Validation reports judge the risk of a model by statistical tests whose error rates the supervisory texts do not state.
We measure the size and power of each test on generators of known truth, under its null and under planted defects, with a Monte Carlo standard error and a Wilson interval for every rate and the exact probability where one exists.
A test holds its size unless its exact rejection probability or the Wilson interval of its rate lies above the level at 5\,\% or 1\,\%.
Of \nTests{} credit-risk tests, \nHold{} hold their size on \nRepsPooled{} null repetitions each and \nNotHold{} do not; the largest excesses are the Jeffreys test, with an exact size of \jeffExactOne\,\% at the 1\,\% level, and the normal test over five years, which rejects \pooledNorm\,\% of correct models at 5\,\% (Wilson interval \pooledNormLo\ to \pooledNormHi\,\%). For \nHoldOneConfig{} of the \nHold{} the null is composite and the size is measured at one configuration.
Cells of \nNullReps{} repetitions resolve \detIntFiveFour{} point and show \nHoldCellDef{} holding, so the count depends on them.
Tests that assume independent defaults reject a correct PD in \binRhoOneZeroExact\,\% of samples at an asset correlation of 10\,\%, the 0.10 rule of the stability index alarms in \psiHundred\,\% of unchanged samples of 100, and the European Central Bank's concentration test cannot reject at 5\,\% with ten grades at an initial coefficient of variation of 0.6.
Of \nPowerTests{} power curves, \nReach{} reach 80\,\% on their ladder; a planted ladder is not a real bank.
